\newpage

\section{Модификации алгоритма Frank--Wolfe}
\label{sec:fw-modifications}

В этой главе рассматриваются модификации алгоритма Frank--Wolfe для задачи поиска равновесного распределения транспортных потоков.
В качестве базовой модели используется модель Бекмана.
Основное внимание уделяется выбору направления шага: классический FW использует только решение текущей линейной подзадачи, тогда как CFW, BFW, FFW, NFW и WFFW строят новое направление с учетом предыдущих направлений.
Такая идея позволяет уменьшить эффект зигзагообразного движения и ускорить уменьшение относительного зазора.

Рассматриваемая линия методов включает CFW и BFW~\cite{mitradjieva2013stiff}, FFW~\cite{fukushima1984modified}, а также последующие расширения~\cite{ignashin2024modifications}:
\begin{enumerate}
    \item Conjugate Frank--Wolfe (CFW);
    \item Bi-conjugate Frank--Wolfe (BFW);
    \item Fukushima Frank--Wolfe (FFW);
    \item N-conjugate Frank--Wolfe (NFW);
    \item Weighted Fukushima Frank--Wolfe (WFFW).
\end{enumerate}
Алгоритмы NFW и WFFW являются развитием BFW и FFW соответственно.
Если BFW использует два предыдущих направления, то NFW использует $N$ предыдущих направлений.
Если FFW использует среднее по нескольким предыдущим направлениям, то WFFW заменяет его экспоненциальным сглаживанием.

\subsection{Модель Бекмана}

Пусть задан ориентированный дорожный граф с множеством ребер $\cE$.
Вектор $f=\{f_e\}_{e\in\cE}$ задает потоки по дорогам, а функция $\tau_e(f_e)$ задает
время движения по ребру $e$ при потоке $f_e$.
В транспортной модели спрос задается не по ребрам, а по корреспонденциям источник--сток.
Пусть $OD$ обозначает множество таких пар:
\[
    OD=\{(i,j)\mid i\in O,\;j\in D\}.
\]
Для каждой пары $w=(i,j)\in OD$ задан спрос $d_w=d_{(i,j)}$, то есть объем потока,
который должен быть перевезен из $i$ в $j$.
Пусть $\cP_w=\cP_{i,j}$ есть множество всех путей из $i$ в $j$, а
$\cP=\bigcup_{w\in OD}\cP_w$ -- множество всех рассматриваемых путей.
Переменная $x_p$ обозначает поток, отправленный по пути $p$.
Тогда для каждой OD-пары допустимые распределения потока по путям образуют симплекс
\[
    X^w=
    \left\{
        x^w \;\middle|\;
        \sum_{p\in\cP_w}x_p=d_w,\;
        x_p\geq 0
    \right\},
\]
а полное множество допустимых путевых потоков имеет декартову структуру
\[
    X=\prod_{w\in OD}X^w.
\]

Связь потоков по путям и потоков по дорогам задается матрицей принадлежности ребер путям:
\[
    f=\Theta x,
    \qquad
    \Theta=\{\delta_{e,p}\}_{e\in\cE,\;p\in\cP},
\]
где $\delta_{e,p}=1$, если путь $p$ проходит по ребру $e$, и $\delta_{e,p}=0$ иначе.
Множество допустимых потоков по ребрам обозначим
\[
    \cF=\{f=\Theta x\mid x\in X\}.
\]

В модели Бекмана равновесное распределение потоков находится как решение задачи минимизации
потенциальной функции
\begin{equation}
    \Psi(f) =
    \sum_{e \in \cE}
    \underbrace{\int_{0}^{f_e} \tau_e(z)\,dz}_{\sigma_e(f_e)}
    \longrightarrow
    \min_{f=\Theta x,\;x\in X}.
    \label{eq:beckmann-primal}
\end{equation}

В качестве функций затрат используются классические BPR-функции:
\begin{equation}
    \tau_e(f_e)=
    \bar t_e
    \left(
        1+\rho\left(\frac{f_e}{\bar f_e}\right)^{1/\mu}
    \right),
    \qquad
    \tau(f)=\{\tau_e(f_e)\}_{e\in\cE}.
    \label{eq:bpr}
\end{equation}
Параметры $\rho,\mu,\bar t_e,\bar f_e$ задаются дорожной сетью.
Градиент функционала $\Psi$ имеет простой вид:
\[
    \nabla \Psi(f)=\tau(f).
\]

\subsection{Оптимальные условия модели Бекмана}
\label{subsec:beckmann-optimality}

Для задачи
\[
    \min_{f=\Theta x,\;x\in X}\Psi(f)
\]
условие оптимальности можно записать как вариационное неравенство
\[
    \inner{\nabla\Psi(f^\ast)}{f-f^\ast}\geq 0,
    \qquad f=\Theta x,\;x\in X.
\]
Так как $\nabla\Psi(f)=\tau(f)$, получаем
\[
    \inner{\tau(f^\ast)}{\Theta x-\Theta x^\ast}\geq 0.
\]
Пусть для некоторой OD-пары $w$ существует используемый путь $p$ и другой путь $q$ с
меньшим временем движения:
\[
    \sum_{e\in q}\tau_e(f^\ast_e)
    <
    \sum_{e\in p}\tau_e(f^\ast_e).
\]
Тогда можно перенести малую величину потока с $p$ на $q$ и получить допустимое направление,
для которого скалярное произведение с $\tau(f^\ast)$ отрицательно.
Это противоречит вариационному неравенству.
Следовательно, в оптимуме все используемые пути являются кратчайшими при равновесных временах.
Так модель Бекмана связывает выпуклую оптимизацию и принцип Вардропа.

\subsection{Линейный оракул как all-or-nothing assignment}
\label{subsec:fw-lmo-aon}

Главная причина, по которой метод Frank--Wolfe удобен для модели Бекмана, состоит в структуре
линейного оракула.
На итерации с текущим потоком $f^k$ нужно решить
\begin{equation}
    s_k^{FW}
    =
    \argmin_{y=\Theta x,\;x\in X}
    \inner{\tau(f^k)}{y}
    =
    \argmin_{x\in X}
    \sum_{w\in OD}\sum_{p\in\cP_w}x_p c_p^k,
    \label{eq:fw-lmo-aon}
\end{equation}
где стоимость пути $p$ при текущих временах равна
\[
    c_p^k
    =
    \sum_{e\in\cE}\delta_{e,p}\tau_e(f_e^k).
\]
Так как $X=\prod_{w\in OD}X^w$, задача~\eqref{eq:fw-lmo-aon} распадается по OD-парам:
\[
    \min_{x^w\in X^w}
    \sum_{p\in\cP_w}x_p c_p^k,
    \qquad
    w=(i,j)\in OD.
\]
Это линейная задача на симплексе.
Поэтому оптимум достигается в вершине: весь спрос $d_w$ назначается на один из путей
минимальной стоимости.
Иными словами,
\[
    p_w^k
    \in
    \argmin_{p\in\cP_w}
    c_p^k
    =
    \argmin_{p\in\cP_w}
    \sum_{e\in\cE}\delta_{e,p}\tau_e(f_e^k),
\]
а решение линейного оракула имеет вид
\[
    \begin{aligned}
        x_{p_w^k}^{k,FW}&=d_w,\\
        x_p^{k,FW}&=0,
        \qquad
        p\in\cP_w\setminus\{p_w^k\},\quad w\in OD,\\
        s_{k,e}^{FW}
        &=
        \sum_{w\in OD}d_w\delta_{e,p_w^k},
        \qquad e\in\cE.
    \end{aligned}
\]
Таким образом, вызов LMO буквально является набором задач поиска кратчайших путей:
для каждой OD-пары $w=(i,j)$ нужно найти кратчайший путь из $i$ в $j$ в графе с длинами
ребер $\tau_e(f_e^k)$.
В реализации не требуется явно перечислять все пути из $\cP_w$; достаточно запустить
алгоритм кратчайших путей по текущим весам ребер.
Если кратчайших путей несколько, можно выбрать любой из них или распределить спрос между
ними без изменения значения линейной подзадачи.

\subsection{Классический алгоритм Frank--Wolfe}

Алгоритм Frank--Wolfe на каждой итерации линеаризует целевой функционал в текущей точке $f^k$:
\[
    \Psi_k(y)=
    \Psi(f^k)+\inner{\nabla\Psi(f^k)}{y-f^k}.
\]
После этого решается линейная задача
\begin{equation}
    s_k^{FW}
    =
    \argmin_{y=\Theta x,\;x\in X}
    \inner{\nabla\Psi(f^k)}{y}
    =
    \argmin_{y=\Theta x,\;x\in X}
    \inner{\tau(f^k)}{y}.
    \label{eq:fw-lmo-traffic}
\end{equation}
Как показано в разделе~\ref{subsec:fw-lmo-aon}, в модели Бекмана эта задача решается как
all-or-nothing assignment по текущим кратчайшим путям.
Направление шага определяется как
\[
    d_k^{FW}=s_k^{FW}-f^k,
\]
а длина шага выбирается одномерной минимизацией:
\[
    \gamma_k=\argmin_{\gamma\in[0,1]}\Psi(f^k+\gamma d_k^{FW}),
    \qquad
    f^{k+1}=f^k+\gamma_k d_k^{FW}.
\]

\begin{algorithm}
\caption{Frank--Wolfe для модели Бекмана}
\begin{algorithmic}
\REQUIRE начальная точка $f^0\in\cF$, максимальное число итераций
\STATE $k\gets 0$
\REPEAT
\STATE $s_k^{FW}\gets \argmin_{s\in\cF}\inner{\nabla\Psi(f^k)}{s}$
\STATE $d_k\gets s_k^{FW}-f^k$
\STATE $\gamma_k\gets\argmin_{\gamma\in[0,1]}\Psi(f^k+\gamma d_k)$
\STATE $f^{k+1}\gets f^k+\gamma_k d_k$
\STATE $k\gets k+1$
\UNTIL{достигнут критерий остановки}
\ENSURE приближение равновесного распределения потоков
\end{algorithmic}
\end{algorithm}

Для оценки качества решения используется зазор Frank--Wolfe
\[
    g_k=
    \max_{s=\Theta x,\;x\in X}
    \inner{\nabla\Psi(f^k)}{f^k-s}.
\]
При липшицевом градиенте известны стандартные оценки сходимости:
\[
    \min_{1\leq m\leq k} g_m
    \leq
    \frac{\min\{2h_0,L\diam(C)^2\}}{\sqrt{k+1}},
    \qquad
    \Psi(f^k)-\Psi(f^\ast)
    \leq
    \frac{2L\diam(C)^2}{k+1}.
\]
Несмотря на эти гарантии, на практике FW часто быстро снижает функционал только на первых итерациях, а затем начинает двигаться зигзагообразно и замедляется.

\subsection{Conjugate Frank--Wolfe}

Идея CFW состоит в построении нового направления, сопряженного к предыдущему направлению относительно гессиана целевой функции.
Пусть $d^{k-1}=s^{k-1}-f^{k-1}$.
Алгоритм ищет новую точку направления в виде выпуклой комбинации
\[
    s_k=\alpha s_k^{FW}+(1-\alpha)s_{k-1}.
\]
Коэффициент $\alpha$ выбирается так, чтобы выполнялось условие сопряженности
\[
    (d^{k-1})^T\nabla^2\Psi(f^k)d^k=0,
    \qquad
    d^k=s_k-f^k.
\]
Для модели Бекмана гессиан диагонален, поэтому умножение на гессиан вычисляется дешево.
Коэффициент $\alpha$ имеет вид
\begin{equation}
    \alpha =
    \frac{
        (f^k-s_{k-1})^T
        \nabla^2\Psi(f^k)
        (s_{k-1}-f^{k-1})
    }{
        (s_k^{FW}-s_{k-1})^T
        \nabla^2\Psi(f^k)
        (s_{k-1}-f^{k-1})
    }.
    \label{eq:cfw-alpha-traffic}
\end{equation}
Чтобы алгоритм не выродился в движение только по старому направлению, коэффициент проецируется на отрезок $[\alpha_0,1]$, где $\alpha_0\in(0,1]$.

\begin{algorithm}
\caption{Conjugate Frank--Wolfe}
\begin{algorithmic}
\REQUIRE $f^0\in\cF$, параметр $\alpha_0\in(0,1]$
\STATE выполнить один шаг FW и получить $f^1$, $s_0$
\STATE $k\gets 1$
\REPEAT
\STATE $s_k^{FW}\gets\argmin_{s\in\cF}\inner{\nabla\Psi(f^k)}{s}$
\STATE вычислить $\alpha$ по формуле~\eqref{eq:cfw-alpha-traffic}
\STATE $\alpha\gets\min\{1,\max\{\alpha_0,\alpha\}\}$
\STATE $s_k\gets\alpha s_k^{FW}+(1-\alpha)s_{k-1}$
\STATE $d_k\gets s_k-f^k$
\STATE $\gamma_k\gets\argmin_{\gamma\in[0,1]}\Psi(f^k+\gamma d_k)$
\STATE $f^{k+1}\gets f^k+\gamma_k d_k$
\STATE $k\gets k+1$
\UNTIL{достигнут критерий остановки}
\end{algorithmic}
\end{algorithm}

\subsection{Fukushima Frank--Wolfe и WFFW}

Алгоритм Fukushima Frank--Wolfe использует усреднение предыдущих направлений.
В общем виде новая точка направления может быть записана как
\begin{equation}
    s_k=
    \alpha_0 s_k^{FW}
    +
    \sum_{j=1}^{N}\alpha_j s_{k-j},
    \qquad
    \sum_{j=0}^{N}\alpha_j=1,
    \qquad
    \alpha_j\geq 0.
    \label{eq:fukushima-general}
\end{equation}
Классический FFW выбирает равномерные веса по нескольким предыдущим направлениям.
Предложенная модификация Weighted Fukushima Frank--Wolfe заменяет равномерное среднее экспоненциальным сглаживанием.
Пусть $\eta\in(0,1)$ задает скорость затухания.
Тогда старые направления получают меньший вес, а более свежие направления сильнее влияют на шаг:
\[
    s_k=
    \alpha_0 s_k^{FW}
    +
    (1-\alpha_0)
    \sum_{j=1}^{N} w_j s_{k-j},
    \qquad
    w_j=
    \frac{(1-\eta)\eta^{j-1}}{\sum_{\ell=1}^{N}(1-\eta)\eta^{\ell-1}}.
\]
Такой выбор сохраняет выпуклость комбинации и одновременно делает алгоритм менее чувствительным к слишком старым направлениям.

\subsection{N-conjugate Frank--Wolfe}

Bi-conjugate Frank--Wolfe использует две предыдущие точки направления.
N-conjugate Frank--Wolfe обобщает эту идею на $N$ предыдущих направлений.
Ищется новая точка
\[
    s_k=
    \alpha_0 s_k^{FW}
    +
    \sum_{j=1}^{N}\alpha_j s_{k-j},
    \qquad
    \alpha_j\geq 0,
    \qquad
    \sum_{j=0}^{N}\alpha_j=1,
\]
такая, что новое направление
\[
    d^k=s_k-f^k
\]
сопряжено к $N$ предыдущим направлениям:
\begin{equation}
    (d^k)^T\nabla^2\Psi(f^k)d^{k-j}=0,
    \qquad
    j=1,\ldots,N.
    \label{eq:nfw-conjugacy}
\end{equation}
В NFW используется предположение, что предыдущие направления уже попарно сопряжены относительно текущего гессиана:
\[
    (d^{k-m})^T\nabla^2\Psi(f^k)d^{k-n}=0,
    \qquad
    m\neq n,\quad m,n=1,\ldots,N.
\]
Тогда коэффициенты могут быть вычислены независимо:
\begin{equation}
    \alpha_j
    =
    -\alpha_0
    \frac{
        (d_k^{FW})^T\nabla^2\Psi(f^k)d^{k-j}
    }{
        (d^{k-j})^T\nabla^2\Psi(f^k)d^{k-j}
    },
    \qquad j=1,\ldots,N,
    \label{eq:nfw-alpha}
\end{equation}
а $\alpha_0$ выбирается так, чтобы сумма весов была равна единице и все веса оставались неотрицательными.

\begin{algorithm}
\caption{N-conjugate Frank--Wolfe}
\begin{algorithmic}
\REQUIRE $f^0\in\cF$, число сопряженных направлений $N$
\STATE выполнить первые $N$ шагов FW или CFW для инициализации
\REPEAT
\STATE $s_k^{FW}\gets\argmin_{s\in\cF}\inner{\nabla\Psi(f^k)}{s}$
\STATE вычислить коэффициенты $\alpha_0,\ldots,\alpha_N$ из условий~\eqref{eq:nfw-conjugacy}
\STATE $s_k\gets\alpha_0s_k^{FW}+\sum_{j=1}^{N}\alpha_js_{k-j}$
\STATE $d_k\gets s_k-f^k$
\STATE $\gamma_k\gets\argmin_{\gamma\in[0,1]}\Psi(f^k+\gamma d_k)$
\STATE $f^{k+1}\gets f^k+\gamma_k d_k$
\UNTIL{достигнут критерий остановки}
\end{algorithmic}
\end{algorithm}

\subsection{Экспериментальные результаты}

Эксперименты проводились на открытых транспортных датасетах из репозитория TransportationNetworks~\cite{bstabler}.
В качестве метрики качества использовался относительный зазор.
Сравнивались FW, CFW, BFW, FFW, WFFW и NFW при различных значениях $N$.

\subsection{Более подробная запись CFW и BFW}

В транспортной постановке CFW рассматривается как первый способ заменить направление
классического FW направлением, которое использует информацию о предыдущем шаге.
Если обозначить обычную FW-точку через $y^k=s_k^{FW}$, то на итерации $k$ строится точка
\[
    \bar y^k=\alpha y^k+(1-\alpha)y^{k-1}.
\]
Далее направление определяется как
\[
    d^k=\bar y^k-f^k.
\]
Коэффициент $\alpha$ получается из условия сопряженности
\[
    (d^k)^T\nabla^2\Psi(f^k)d^{k-1}=0.
\]
Подставляя $d^k=\alpha(y^k-f^k)+(1-\alpha)(y^{k-1}-f^k)$, получаем
\[
    \alpha
    =
    \frac{(f^k-y^{k-1})^T\nabla^2\Psi(f^k)d^{k-1}}
    {(y^k-y^{k-1})^T\nabla^2\Psi(f^k)d^{k-1}}.
\]
После вычисления коэффициент ограничивается снизу параметром $\alpha_{\min}$:
\[
    \alpha\leftarrow \max\{\alpha_{\min},\min\{1,\alpha\}\}.
\]
Это техническое ограничение важно: оно не позволяет методу полностью игнорировать новое решение
линейного оракула и тем самым сохраняет связь с классическим FW.

Bi-conjugate Frank--Wolfe использует уже две предыдущие точки.
Пусть нужно построить
\[
    \bar y^k=\alpha_0y^k+\alpha_1y^{k-1}+\alpha_2y^{k-2},
    \qquad
    \alpha_0+\alpha_1+\alpha_2=1.
\]
Условия сопряженности имеют вид
\[
    (d^k)^T\nabla^2\Psi(f^k)d^{k-1}=0,
    \qquad
    (d^k)^T\nabla^2\Psi(f^k)d^{k-2}=0.
\]
Эта идея используется как промежуточный шаг к NFW: BFW показывает,
что память о двух направлениях может быть лучше памяти об одном направлении,
но жестко фиксированное число направлений не позволяет проверить, где находится
практически оптимальный уровень памяти алгоритма.

\begin{algorithm}
\caption{Bi-conjugate Frank--Wolfe}
\begin{algorithmic}
\REQUIRE допустимая точка $f^0$, нижняя граница $\alpha_{\min}$, число итераций
\STATE выполнить две начальные итерации FW/CFW и сохранить $y^0,y^1$
\FOR{$k=2,3,\ldots$}
\STATE $y^k\gets\argmin_{y\in\cF}\inner{\nabla\Psi(f^k)}{y}$
\STATE найти $\alpha_0,\alpha_1,\alpha_2$ из условий сопряженности к $d^{k-1}$ и $d^{k-2}$
\STATE спроецировать коэффициенты на допустимый симплекс с учетом $\alpha_0\geq\alpha_{\min}$
\STATE $\bar y^k\gets\alpha_0y^k+\alpha_1y^{k-1}+\alpha_2y^{k-2}$
\STATE $d^k\gets\bar y^k-f^k$
\STATE $\gamma_k\gets\argmin_{\gamma\in[0,1]}\Psi(f^k+\gamma d^k)$
\STATE $f^{k+1}\gets f^k+\gamma_kd^k$
\ENDFOR
\end{algorithmic}
\end{algorithm}

\subsection{Вывод коэффициентов NFW}

Пусть на итерации $k$ доступны $N$ предыдущих направлений
\[
    d^{k-1},d^{k-2},\ldots,d^{k-N}.
\]
NFW ищет новую точку в виде
\[
    \bar y^k=\alpha_0y^k+\sum_{i=1}^{N}\alpha_i y^{k-i},
    \qquad
    \sum_{i=0}^{N}\alpha_i=1.
\]
Тогда
\[
    d^k=\bar y^k-f^k
    =
    \alpha_0(y^k-f^k)
    +
    \sum_{i=1}^{N}\alpha_i(y^{k-i}-f^k).
\]
Для каждого $j=1,\ldots,N$ требуется
\[
    (d^k)^T H_k d^{k-j}=0,
    \qquad
    H_k=\nabla^2\Psi(f^k).
\]
Если предыдущие направления считаются попарно сопряженными относительно текущего $H_k$,
то есть
\[
    (d^{k-i})^T H_k d^{k-j}=0,\qquad i\neq j,
\]
коэффициенты можно вычислять последовательно.
Это предположение техническое, но оно согласуется с мотивацией метода:
если на предыдущих итерациях направления были сопряжены относительно близких гессианов,
а $H_k$ меняется плавно, то старые направления остаются почти сопряженными и относительно
нового гессиана.

Выведем формулы подробнее.
Из обновления
\[
    f^k=(1-\gamma_{k-1})f^{k-1}+\gamma_{k-1}y^{k-1},
    \qquad
    d^{k-1}=y^{k-1}-f^{k-1},
\]
получаем
\[
    y^{k-1}-f^k=(1-\gamma_{k-1})d^{k-1}.
\]
Обозначим
\[
    \bar d^{k-i}=(1-\gamma_{k-i})d^{k-i}.
\]
Для $m\geq 1$ используем телескопическую сумму:
\[
    y^{k-m}-f^k
    =
    \sum_{i=0}^{m-2}
    \left(y^{k-m+i}-y^{k-m+i+1}\right)
    +
    y^{k-1}-f^k .
\]
Разность соседних точек линейного оракула переписывается через прошлые направления:
\[
    y^{k-m+i}-y^{k-m+i+1}
    =
    \bar d^{k-m+i}-d^{k-m+i+1}.
\]
Следовательно,
\begin{equation}
    d^k
    =
    \alpha_0d_k^{FW}
    +
    \sum_{m=1}^{N}
    \alpha_m
    \left(
        \sum_{i=0}^{m-1}\bar d^{k-m+i}
        -
        \sum_{i=0}^{m-2}d^{k-m+i+1}
    \right).
    \label{eq:nfw-direction-expanded}
\end{equation}
Теперь подставим~\eqref{eq:nfw-direction-expanded} в условия сопряженности
\[
    (d^{k-m})^T H_k d^k=0,
    \qquad m=1,\ldots,N.
\]
Введем обозначения
\[
    A_{k-m}=(d^{k-m})^T H_k d_k^{FW},
    \qquad
    B_{k-m}=(d^{k-m})^T H_k d^{k-m}.
\]
После зануления смешанных слагаемых по предположению о попарной сопряженности остается
\begin{equation}
    \alpha_0A_{k-m}
    +
    (1-\gamma_{k-m})B_{k-m}\alpha_m
    -
    \gamma_{k-m}B_{k-m}
    \sum_{i=m+1}^{N}\alpha_i
    =
    0.
    \label{eq:nfw-alpha-system}
\end{equation}
Положим $\alpha_m=\beta_m\alpha_0$.
Тогда коэффициенты $\beta_m$ считаются с конца к началу:
\begin{equation}
    \beta_N
    =
    -\frac{A_{k-N}}{B_{k-N}(1-\gamma_{k-N})},
    \label{eq:nfw-beta-last}
\end{equation}
\begin{equation}
    \beta_m
    =
    -\frac{A_{k-m}}{B_{k-m}(1-\gamma_{k-m})}
    +
    \frac{\gamma_{k-m}}{1-\gamma_{k-m}}
    \sum_{i=m+1}^{N}\beta_i,
    \qquad m=N-1,\ldots,1.
    \label{eq:nfw-beta-recursion}
\end{equation}
Нормировка выпуклой комбинации дает
\begin{equation}
    \alpha_0=
    \frac{1}{1+\sum_{i=1}^{N}\beta_i},
    \qquad
    \alpha_m=\beta_m\alpha_0.
    \label{eq:nfw-alpha-normalization}
\end{equation}
Если найденные коэффициенты нарушают неотрицательность, алгоритм возвращается к более простой
комбинации: уменьшает число используемых направлений или фактически делает шаг CFW/FW.
Именно поэтому NFW не является отдельным дорогим методом: он является надстройкой над тем же
линейным оракулом, а дополнительные вычисления состоят из нескольких скалярных произведений с
диагональным гессианом.

Отдельная эвристика связана с шагом $\gamma_k=1$.
Если такой шаг выполняется, новая точка полностью задается текущим направлением, а накопленная
информация о прошлых направлениях теряет смысл.
Поэтому в реализации метода используется параметр $\gamma_{\max}$:
если $\gamma_k\geq\gamma_{\max}$, счетчик доступных сопряженных направлений сбрасывается,
и накопление памяти для NFW начинается заново.

\subsection{Fukushima Frank--Wolfe как сглаживание направлений}

В алгоритме Fukushima Frank--Wolfe новое направление не выводится из условий сопряженности.
Оно строится как выпуклое среднее по текущему FW-направлению и нескольким предыдущим точкам
линейного оракула.
Если хранится $N$ прошлых точек, то
\[
    \bar y^k
    =
    \lambda y^k
    +
    (1-\lambda)
    \frac{1}{N}\sum_{i=1}^{N}y^{k-i}.
\]
После этого
\[
    d^k=\bar y^k-f^k.
\]
Такой подход проще CFW/NFW: он не использует гессиан и не пытается сделать направления
сопряженными.
Его смысл состоит в подавлении резких колебаний, возникающих из-за того, что линейный оракул
часто выбирает сильно различающиеся кратчайшие пути на соседних итерациях.

Weighted FFW меняет только способ усреднения:
\[
    \bar y^k
    =
    \lambda y^k
    +
    (1-\lambda)
    \sum_{i=1}^{N}w_i y^{k-i},
    \qquad
    \sum_{i=1}^{N}w_i=1.
\]
Веса выбираются геометрическими:
\[
    w_i=\frac{q^{i-1}}{\sum_{\ell=1}^{N}q^{\ell-1}},
    \qquad q\in(0,1).
\]
При малом $q$ метод почти полностью ориентируется на последнее сохраненное направление.
При $q$ около единицы WFFW становится близким к обычному FFW.
Такой параметр позволяет управлять тем, насколько быстро алгоритм забывает устаревшую информацию.

\begin{algorithm}
\caption{Weighted Fukushima Frank--Wolfe}
\begin{algorithmic}
\REQUIRE $f^0\in\cF$, глубина памяти $N$, параметры $\lambda,q\in(0,1)$
\FOR{$k=0,1,2,\ldots$}
\STATE $y^k\gets\argmin_{y\in\cF}\inner{\nabla\Psi(f^k)}{y}$
\STATE выбрать доступные прошлые точки $y^{k-1},\ldots,y^{k-N}$
\STATE $w_i\gets q^{i-1}/\sum_{\ell=1}^{N}q^{\ell-1}$
\STATE $\bar y^k\gets\lambda y^k+(1-\lambda)\sum_{i=1}^{N}w_i y^{k-i}$
\STATE $d^k\gets\bar y^k-f^k$
\STATE $\gamma_k\gets\argmin_{\gamma\in[0,1]}\Psi(f^k+\gamma d^k)$
\STATE $f^{k+1}\gets f^k+\gamma_kd^k$
\ENDFOR
\end{algorithmic}
\end{algorithm}

\subsection{Относительный зазор и критерий остановки}

В экспериментах основной метрикой является относительный зазор, используемый в работах по
ускоренным FW-методам для traffic assignment~\cite{mitradjieva2013stiff,arrache2008accelerating}.
В отличие от простого значения целевой функции, эта величина сравнивает текущее верхнее
приближение оптимального значения с накопленной нижней оценкой минимума.
Поэтому она показывает не только факт убывания $\Psi$, но и то, насколько близко текущее
решение к оптимальному значению в относительном масштабе.

Пусть $f^k$ -- текущий поток, $s_k^{FW}$ -- решение линейного оракула, а $f^{k+1}$ -- точка
после линейного поиска.
Верхняя оценка оптимального значения берется как значение целевой функции в новой допустимой
точке:
\begin{equation}
    \operatorname{UBD}_k=\Psi(f^{k+1}).
    \label{eq:ubd-relgap}
\end{equation}
Так как $f^{k+1}\in\cF$, всегда выполнено
$\Psi(f^\ast)\leq \operatorname{UBD}_k$.

Нижняя оценка строится из линеаризации выпуклой функции $\Psi$ в точке $f^k$:
\begin{equation}
    \operatorname{LB}_k
    =
    \min_{y=\Theta x,\;x\in X}
    \left\{
        \Psi(f^k)+\inner{\nabla\Psi(f^k)}{y-f^k}
    \right\}.
    \label{eq:lb-relgap}
\end{equation}
По выпуклости для любой допустимой точки $y$ справедливо
\[
    \Psi(y)\geq \Psi(f^k)+\inner{\nabla\Psi(f^k)}{y-f^k}.
\]
В частности, для оптимального потока $f^\ast$ получаем
\[
    \Psi(f^\ast)
    \geq
    \Psi(f^k)+\inner{\nabla\Psi(f^k)}{f^\ast-f^k}
    \geq
    \operatorname{LB}_k.
\]
Следовательно, $\operatorname{LB}_k$ действительно является нижней оценкой оптимального
значения.

Эту нижнюю оценку можно выразить через зазор Frank--Wolfe.
Обозначим
\begin{equation}
    g_k
    =
    \inner{\nabla\Psi(f^k)}{f^k-s_k^{FW}}
    =
    \inner{\tau(f^k)}{f^k-s_k^{FW}}.
    \label{eq:fw-gap-relgap-section}
\end{equation}
Так как $s_k^{FW}$ минимизирует линейную форму $\inner{\nabla\Psi(f^k)}{y}$ на допустимом
множестве, из~\eqref{eq:lb-relgap} следует
\begin{equation}
    \operatorname{LB}_k
    =
    \Psi(f^k)+\inner{\nabla\Psi(f^k)}{s_k^{FW}-f^k}
    =
    \Psi(f^k)-g_k.
    \label{eq:lb-fw-gap}
\end{equation}
В транспортной постановке $g_k$ имеет понятную интерпретацию: он измеряет, насколько можно
уменьшить линейную аппроксимацию целевой функции, если переназначить весь спрос на текущие
кратчайшие пути.

Поскольку отдельные нижние оценки могут колебаться от итерации к итерации, в относительном
зазоре используется лучшая накопленная нижняя оценка:
\begin{equation}
    \operatorname{BLB}_k
    =
    \max_{1\leq m\leq k}\operatorname{LB}_m.
    \label{eq:blb-relgap}
\end{equation}
Из $\operatorname{LB}_m\leq\Psi(f^\ast)$ для всех $m$ следует
$\operatorname{BLB}_k\leq\Psi(f^\ast)$.
Итоговая метрика задается формулой
\begin{equation}
    \operatorname{RG}_k
    =
    \frac{\operatorname{UBD}_k-\operatorname{BLB}_k}
         {\operatorname{BLB}_k}.
    \label{eq:rg-relgap}
\end{equation}
Если $\operatorname{BLB}_k>0$, что выполняется для стандартных транспортных задач с
положительными временами движения, то
\begin{equation}
    \frac{\Psi(f^{k+1})-\Psi(f^\ast)}{\Psi(f^\ast)}
    \leq
    \frac{\Psi(f^{k+1})-\operatorname{BLB}_k}{\operatorname{BLB}_k}
    =
    \operatorname{RG}_k.
    \label{eq:rg-upper-bound}
\end{equation}
Таким образом, $\operatorname{RG}_k$ является вычислимой верхней оценкой относительной
ошибки по целевой функции.

Также эта величина не хуже грубой относительной оценки, построенной напрямую по FW-зазору.
Действительно, при точном линейном поиске $\Psi(f^{k+1})\leq\Psi(f^k)$, а
$\operatorname{BLB}_k\geq\operatorname{LB}_k$.
Поэтому
\begin{equation}
    \operatorname{RG}_k
    =
    \frac{\Psi(f^{k+1})-\operatorname{BLB}_k}{\operatorname{BLB}_k}
    \leq
    \frac{\Psi(f^k)-\operatorname{LB}_k}{\operatorname{BLB}_k}
    =
    \frac{g_k}{\operatorname{BLB}_k}.
    \label{eq:rg-vs-fw-gap}
\end{equation}
Именно поэтому в сравнении модификаций FW используется $\operatorname{RG}_k$: он сохраняет
связь с зазором Frank--Wolfe, но дополнительно учитывает лучшую накопленную нижнюю оценку
оптимального значения.
Алгоритм считается более быстрым, если достигает одного и того же уровня
$\operatorname{RG}_k$ за меньшее число итераций или за меньшее время.

Для всех модификаций схема эксперимента одинакова:
\begin{enumerate}
    \item построить начальный допустимый поток назначением спроса на кратчайшие пути при свободных временах;
    \item на каждой итерации пересчитать времена $\tau(f^k)$;
    \item решить линейный оракул как набор задач кратчайших путей;
    \item построить направление конкретной модификации;
    \item выполнить точный линейный поиск по $\gamma\in[0,1]$;
    \item вычислить $\operatorname{LB}_k$, обновить $\operatorname{BLB}_k$ и сохранить
    $\operatorname{RG}_k$ вместе со значением целевой функции.
\end{enumerate}
Такая унификация важна: различия в результатах объясняются именно выбором направления,
а не разными правилами шага или разными линейными оракулами.

\subsection{Серии экспериментов}

Эксперименты проводились на транспортных сетях Anaheim, SiouxFalls, Berlin-Tiergarten,
Terrassa-Asymmetric, Chicago-Sketch, Berlin-Mitte-Center, Berlin-Friedrichshain, Barcelona
и Berlin-Mitte-Prenzlauerberg-Friedrichshain-Center.
Во всех сериях сравнивается зависимость относительного зазора от времени работы.
Это важно: методы используют один и тот же линейный оракул кратчайших путей и один и тот же
линейный поиск, поэтому различия на графиках отражают именно выбор направления.

Первая серия сравнивает FW, CFW, BFW, WFFW и $3$-conjugate FW.
Главный вопрос здесь состоит в том, дает ли третье сопряженное направление выигрыш по сравнению
с уже известным BFW.
На рис.~\ref{fig:fw-mod-cfw-0}--\ref{fig:fw-mod-cfw-8} видно, что на многих сетях
$3$-conjugate FW уменьшает относительный зазор быстрее BFW, CFW и классического FW.
Это означает, что эффект памяти не исчерпывается двумя направлениями.

\begin{figure}[H]
\centering
\includegraphics[width=0.98\textwidth]{fw_modifications/0_CFW.pdf}
\caption{Сходимость относительного зазора для NFW, CFW, BFW, WFFW и FW на сетях Chicago-Sketch, Berlin-Friedrichshain и Anaheim.}
\label{fig:fw-mod-cfw-0}
\end{figure}

\begin{figure}[H]
\centering
\includegraphics[width=0.98\textwidth]{fw_modifications/3_CFW.pdf}
\caption{Сходимость относительного зазора для NFW, CFW, BFW, WFFW и FW на сетях Berlin-Mitte-Center, Berlin-Tiergarten и Barcelona.}
\label{fig:fw-mod-cfw-3}
\end{figure}

\begin{figure}[H]
\centering
\includegraphics[width=0.98\textwidth]{fw_modifications/8_CFW.pdf}
\caption{Сходимость относительного зазора для NFW, CFW, BFW, WFFW и FW на сетях Berlin-Mitte-Prenzlauerberg-Friedrichshain-Center, SiouxFalls и Terrassa-Asymmetric.}
\label{fig:fw-mod-cfw-8}
\end{figure}

Вторая серия посвящена сравнению Fukushima Frank--Wolfe и Weighted Fukushima Frank--Wolfe.
Здесь проверяется, насколько простое сглаживание направлений может конкурировать с
сопряженными методами.
Равномерное усреднение FFW стабилизирует траекторию, но может запаздывать, поскольку старые
направления удерживают алгоритм в области, которая уже не соответствует текущим кратчайшим
путям.
WFFW исправляет эту особенность с помощью геометрического забывания.
Результаты приведены на рис.~\ref{fig:fw-mod-wffw-0}--\ref{fig:fw-mod-wffw-2}.

\begin{figure}[H]
\centering
\includegraphics[width=0.98\textwidth]{fw_modifications/0_WFW_weighted.pdf}
\caption{Сравнение FFW, WFFW и CFW по относительному зазору на сетях Anaheim и Berlin-Friedrichshain.}
\label{fig:fw-mod-wffw-0}
\end{figure}

\begin{figure}[H]
\centering
\includegraphics[width=0.98\textwidth]{fw_modifications/2_WFW_weighted.pdf}
\caption{Сравнение FFW, WFFW и CFW по относительному зазору на сетях SiouxFalls и Terrassa-Asymmetric.}
\label{fig:fw-mod-wffw-2}
\end{figure}

Третья серия меняет глубину памяти NFW.
Сравниваются варианты с разным числом предыдущих направлений.
При $N=2$ метод совпадает по смыслу с BFW, а при больших $N$ начинает использовать более
длинную историю направлений.
На рис.~\ref{fig:fw-mod-nfw-0}--\ref{fig:fw-mod-nfw-2} видно, что увеличение глубины памяти
может давать дополнительное ускорение, но эффект зависит от сети: слишком старые направления
иногда описывают уже неактуальную локальную геометрию задачи.

\begin{figure}[H]
\centering
\includegraphics[width=0.98\textwidth]{fw_modifications/0_NFW.pdf}
\caption{Сравнение NFW с разной глубиной памяти и BFW на сетях Anaheim и Berlin-Friedrichshain.}
\label{fig:fw-mod-nfw-0}
\end{figure}

\begin{figure}[H]
\centering
\includegraphics[width=0.98\textwidth]{fw_modifications/2_NFW.pdf}
\caption{Сравнение NFW с разной глубиной памяти и BFW на сетях SiouxFalls и Terrassa-Asymmetric.}
\label{fig:fw-mod-nfw-2}
\end{figure}

Четвертая серия проверяет параметры WFFW.
Если вес текущего FW-направления слишком мал, метод двигается инерционно и может медленно
реагировать на изменение кратчайших путей.
Если этот вес слишком велик, метод почти совпадает с FW и теряет смысл сглаживания.
На рис.~\ref{fig:fw-mod-wffw-params} показано, что нужен баланс между новым FW-направлением
и памятью о прошлых направлениях.

\begin{figure}[H]
\centering
\includegraphics[width=0.98\textwidth]{fw_modifications/0_WFW.pdf}
\caption{Сходимость WFFW при различных параметрах геометрического сглаживания на сети Terrassa-Asymmetric.}
\label{fig:fw-mod-wffw-params}
\end{figure}

Основные наблюдения следующие.
Во-первых, методы, использующие предыдущие направления, в большинстве экспериментов сходятся
быстрее классического FW.
Во-вторых, использование более чем двух сопряженных направлений может давать дополнительное
ускорение: вариант $3$-conjugate FW часто оказывается одним из наиболее эффективных.
В-третьих, WFFW обычно улучшает FFW и CFW за счет управляемого забывания старых направлений,
хотя в ряде экспериментов уступает NFW.
Таким образом, использование памяти о предыдущих направлениях является устойчивым способом
ускорить FW без изменения линейного оракула.

\subsection{Одномерный поиск для BPR-функций}

Длина шага выбирается точной минимизацией
\[
    \gamma_k=\argmin_{\gamma\in[0,1]}\Psi(f^k+\gamma d^k).
\]
Для BPR-функции
\[
    \tau_e(f_e)=\bar t_e\left(1+\rho\left(\frac{f_e}{\bar f_e}\right)^{1/\mu}\right)
\]
производная одномерной функции равна
\[
    \frac{d}{d\gamma}\Psi(f^k+\gamma d^k)
    =
    \sum_{e\in\cE}
    \tau_e(f_e^k+\gamma d_e^k)d_e^k.
\]
Поэтому задача поиска $\gamma$ сводится к нахождению нуля монотонной функции на $[0,1]$.
Если производная в нуле неотрицательна, оптимум находится при $\gamma=0$.
Если производная в единице неположительна, оптимум находится при $\gamma=1$.
В остальных случаях используется одномерный численный метод.
Эта процедура одинакова для FW, CFW, BFW, FFW, WFFW и NFW; различается только направление
$d^k$.

\subsection{Итоги раздела}

Рассмотренное направление развития движется от уже известных CFW/BFW к двум новым вариантам.
NFW проверяет гипотезу, что память о большем числе прошлых направлений может усилить
эффект сопряженности.
WFFW проверяет другую гипотезу: может быть, часть ускорения объясняется не строгой
сопряженностью, а более устойчивым сглаживанием направлений.
Эксперименты показывают, что обе гипотезы полезны.
Сопряженные методы дают наиболее сильный эффект, но взвешенное сглаживание также улучшает
поведение классического FW.


\subsection{Выводы по модификациям}

Рассмотренные модификации сохраняют ключевое преимущество FW: каждая итерация требует решения той же линейной подзадачи, которая в транспортной модели сводится к поиску кратчайших путей.
Дополнительные вычисления связаны с линейным поиском и комбинированием предыдущих направлений.
Для модели Бекмана эти операции существенно дешевле линейного оракула, поскольку гессиан диагонален.

Эксперименты показывают, что направление шага является главным объектом для улучшения FW.
Наиболее перспективным оказывается семейство сопряженных методов, особенно NFW при малых значениях $N$.
Эта идея становится отправной точкой для следующей главы, где CFW и NFW переносятся из транспортной постановки в задачи машинного обучения.
